\section*{Chapitre \ref{ch:pl:the-research-environment} --- L'environnement de recherche}

\begin{solution}{exo:pl:the-research-environment:1}
La quatrième cellule de code, dont le compteur vaut 3, est hors d'ordre (la cellule au-dessus d'elle a le compteur 5)~; le compteur 4 manque (une réexécution ou une cellule supprimée). La
troisième cellule s'est exécutée après la quatrième, et autre chose s'est exécuté entre les deux.
\end{solution}

\begin{solution}{exo:pl:the-research-environment:2}
Aucune cellule du notebook n'affecte \texttt{exclude}, ni au-dessus ni en dessous~: le nom a été lié par une cellule qui a été supprimée. «~Utilisé avant d'être défini~» sert pour un
nom qu'une cellule plus bas affecte --- une réexécution depuis le haut échoue là aussi, mais la définition figure toujours dans le document.
\end{solution}

\begin{solution}{exo:pl:the-research-environment:3}
Les cinq pays exclus par la cellule supprimée, et les années antérieures à 1950, qui étaient encore dans les données lorsque le résultat a été calculé parce que la
cellule de nettoyage s'est exécutée après. Avec l'exclusion et toutes les années, la bibliothèque donne 1.96~; avec tous les pays à partir de 1950, 2.44.
\end{solution}

\begin{solution}{exo:pl:the-research-environment:4}
Oui. Une cellule modifiée après son exécution affiche sa nouvelle source avec l'ancienne sortie~; une cellule qui ajoute à une liste, exécutée une fois avec des compteurs consécutifs et
puis réexécutée après un redémarrage qui a réutilisé des numéros de compteur, ne laisse rien paraître dans les compteurs~; un nom défini à l'invite d'une console attachée
ne laisse aucune trace. Seule la réexécution est concluante.
\end{solution}

\begin{solution}{exo:pl:the-research-environment:5}
Les deux répondent à une question~; ce sont des questions différentes --- l'expérience du pays type face à celle du couple pays-année type --- et la
différence (2.44 contre 2.90) est plus grande que bien des effets publiés. Une fonction qui masque le choix publie une réponse comme si c'était la
seule~; en argument, il figure dans l'appel et dans le rapport. Présenter celle qui correspond à la question, et l'autre à côté.
\end{solution}

\begin{solution}{exo:pl:the-research-environment:6}
Le verrouillage de l'environnement suffit lorsque le résultat ne dépend que de l'interpréteur et des paquets Python qu'il fixe, sur un système d'exploitation connu~; un conteneur est
nécessaire lorsqu'il dépend de bibliothèques système, de compilateurs ou d'outils situés hors du verrouillage (\omterm{def:pl:native-extensions:extension}{extensions natives}, clients de bases de données), ou doit tourner sur des machines
que l'on ne contrôle pas.
\end{solution}

\begin{solution}{exo:pl:the-research-environment:7}
Sur le notebook corrigé, la vérification n'a rien à comparer~: la correction a effacé les sorties enregistrées, comme il se doit, puisqu'elles provenaient d'un historique qui ne
existe plus. Sur le notebook piégé, la réexécution s'arrête à la quatrième cellule de code, si bien que le 1.96 enregistré dans la cellule de résultat n'est jamais comparé à rien.
C'est plus difficile que de comparer la dernière sortie, car la sortie de la réexécution doit être rattachée aux cellules (des marqueurs imprimés entre elles), certaines sorties
diffèrent légitimement (heures, adresses mémoire, derniers chiffres des sommes en virgule flottante), et les sorties riches, comme les graphiques, demandent une comparaison propre.
\end{solution}

\begin{solution}{exo:pl:the-research-environment:8}
Le contrôle de version enregistre le document, avec ses sorties enregistrées~; il n'enregistre pas l'état du noyau qui les a produites, et le relecteur lit le
document. Reproductible signifie réexécutable~: depuis une copie de travail propre, dans l'environnement verrouillé, sur l'instantané des données, depuis le haut, avec le même
résultat. Ni le commit ni l'approbation n'ont fait cela.
\end{solution}

\begin{solution}{pb:pl:the-research-environment:1}
\begin{enumerate}
\item Sa source, ses sorties enregistrées et le compteur d'exécution du noyau à sa dernière exécution (un entier, ou null).
\item Exécution~: chargement, sélection au-dessus de 90~\%, l'exclusion supprimée, la sélection des pays, le résultat, puis le nettoyage~; affiché~: chargement, nettoyage,
sélection, sélection des pays, résultat.
\item 1.96.
\item Une cellule exécutée hors d'ordre, une cellule réexécutée, une cellule modifiée ou supprimée après son exécution (et les noms définis à l'invite).
\item Sur 863\,878 exécutions tentées de notebooks valides, 24.11~\% se sont déroulées sans erreur et 4.03~\% ont reproduit leurs résultats.
\item Hors d'ordre (la cellule 3 s'est exécutée avant la cellule 2), un compteur sauté (3), et un nom défini nulle part (\texttt{exclude})~; à partir du seul JSON.
\item Elle échoue~: \texttt{NameError} sur \texttt{exclude}.
\item Le compteur sauté et le nom non défini correspondent à l'exclusion supprimée~; le compteur hors d'ordre indique que les données affichées ne sont pas celles utilisées par le
résultat.
\item Les modifications après une exécution, les réexécutions qui laissent des compteurs consécutifs, et tout ce qui est fait en dehors du notebook.
\item Pour que rien de la session de l'auteur --- noms, imports, répertoire de travail --- ne s'infiltre dans la vérification.
\item Le seuil, la première année, les pays exclus et la pondération.
\item 1.96 (cinq pays exclus, toutes les années), 2.44 (tous les pays à partir de 1950, par pays), 2.90 (tous les pays à partir de 1950, par couple pays-année).
\item 204 sur 455.
\item La liste de contrôle~: aucun \omterm{def:pl:the-research-environment:hidden}{état caché}, réexécution réussie et reproductible, tests, environnement confronté au verrouillage, données confrontées à l'instantané.
\item La bibliothèque et chaque appelant évoluent ensemble et sont testés ensemble.
\item \textbf{Résultat nommé.} Le notebook affiche 1.96~; réexécuté depuis le haut, il échoue sur un nom défini seulement par une cellule supprimée~; l'audit trouve cette
cellule (un compteur sauté et un nom défini nulle part) et une cellule de nettoyage exécutée après le résultat~; corrigé, le calcul donne 2.44, et la bibliothèque
promue la reproduit depuis un répertoire propre, avec l'instantané des données et le verrouillage de l'environnement vérifiés.
\item Le calcul promu, une réexécution depuis une copie de travail propre qui l'affiche, le verrouillage et l'instantané des données, et les choix énoncés dans
l'appel.
\item Lorsque le résultat dépend de quoi que ce soit hors du verrouillage, ou doit tourner sur des machines que je ne contrôle pas.
\item La question, les choix (seuil, années, pays, pondération), la commande qui recalcule le résultat, les empreintes du code et des données, et
les alternatives essayées.
\item Que n'importe qui puisse le recalculer depuis le haut, à partir d'une copie de travail propre, et obtienne le même nombre.
\end{enumerate}
\end{solution}

\begin{solution}{iq:pl:the-research-environment:1}
Le réexécuter depuis le haut dans un noyau neuf~; lire ses compteurs d'exécution pour repérer les sauts et les exécutions hors d'ordre~; chercher les noms définis nulle part~; vérifier la
version des données et des paquets par rapport à ce qu'avait utilisé l'auteur.
\iqlookfor{Réexécuter d'abord, puis \omterm{def:pl:the-research-environment:hidden}{état caché}, données et environnement.}
\end{solution}

\begin{solution}{iq:pl:the-research-environment:2}
Dès qu'il est utilisé deux fois, qu'il alimente ce dont d'autres dépendent, ou que ses choix comptent~: il faut alors des arguments pour ces choix, des tests et une \omterm{def:pl:the-research-environment:review}{revue de code}.
\iqlookfor{Réutilisation et conséquences, et non la taille.}
\end{solution}

\begin{solution}{iq:pl:the-research-environment:3}
L'ensemble de l'environnement sous les paquets Python --- système d'exploitation, bibliothèques système, compilateurs --- figé et portable. Cela a un coût en temps de construction,
en stockage d'images et en maintenance, et une \omterm{def:pl:the-research-environment:container}{image de conteneur} peut toujours cacher ce qu'elle contient, à moins d'être construite à partir d'un verrouillage.
\iqlookfor{Figer plutôt que décrire~; construire à partir d'un verrouillage.}
\end{solution}

\begin{solution}{iq:pl:the-research-environment:4}
Produire chaque nombre par une commande du dépôt qui s'exécute depuis une copie de travail propre, sur un instantané de données, dans un environnement verrouillé, et consigner dans le
rapport le commit, l'empreinte de l'instantané et la commande~; tester en intégration continue que les commandes reproduisent encore les nombres.
\iqlookfor{Les nombres comme sorties de commandes versionnées.}
\end{solution}

\begin{solution}{iq:pl:the-research-environment:5}
Des notebooks pour l'exploration, une \omterm{def:pl:the-research-environment:library}{bibliothèque de recherche} commune et testée dans un même dépôt, avec \omterm{def:pl:the-research-environment:review}{revue de code}, des environnements issus d'un verrouillage (conteneurs là où c'est nécessaire),
des données lues à travers les instantanés de la plateforme, un contrôle d'audit et de réexécution avant tout partage de résultat, et un journal de recherche avec les commandes qui
reproduisent chaque résultat.
\iqlookfor{Exploration libre, résultats filtrés.}
\end{solution}
